\documentclass[../DoAn.tex]{subfiles}
\begin{document}

\begin{center}
    \Large{\textbf{TÓM TẮT NỘI DUNG ĐỒ ÁN}}\\
\end{center}
\vspace{1cm}

Trong những năm gần đây, ngành công nghiệp trò chơi điện tử chứng kiến sự phát triển mạnh mẽ của thể loại Metroidvania kết hợp hành động đi cảnh (Action Platformer), tiêu biểu là các tác phẩm như Hollow Knight hay Never Grave. Dù đạt nhiều thành công, việc phát triển các trò chơi thuộc thể loại này vẫn đối mặt với thách thức lớn trong việc cân bằng giữa tính khám phá thế giới mở và duy trì trải nghiệm chiến đấu thử thách. Các hướng tiếp cận hiện tại thường tập trung vào thiết kế màn chơi tuyến tính hoặc đơn giản hóa hệ thống trí tuệ nhân tạo (AI) của kẻ địch, dẫn đến hạn chế về chiều sâu trải nghiệm và giảm giá trị chơi lại của sản phẩm. Để giải quyết vấn đề này, đồ án lựa chọn hướng tiếp cận nghiên cứu và tối ưu hóa quy trình xây dựng game Metroidvania trên nền tảng Unity, nhằm tạo ra một hệ thống cơ chế vận hành phức tạp nhưng vẫn đảm bảo hiệu năng ổn định. 

Giải pháp tổng quan được triển khai thông qua dự án trò chơi "Dream Witch". Tựa game được phát triển dựa trên việc tích hợp toàn diện cấu trúc thế giới mở liên thông, hệ thống chiến đấu chặt chẽ và các cơ chế xây dựng nhân vật có chiều sâu. Điểm nổi bật về mặt kỹ thuật và cũng là đóng góp chính của đồ án bao gồm việc thiết kế thành công hệ thống nhập hồn Enemy, nâng cao độ đa dạng trải nghiệm của người chơi. Kết quả sau cùng của đồ án là một sản phẩm trò chơi hoàn chỉnh, vận hành ổn định trên môi trường Unity, chứng minh tính khả thi của các giải pháp kỹ thuật đã đề xuất và mở ra hướng tối ưu hóa hiệu quả cho quy trình phát triển game độc lập.

\begin{flushright}
Sinh viên thực hiện\\
\vspace{1.5cm} % Khoảng trống để ký tên
\begin{tabular}{@{}c@{}}
\textit{(Ký và ghi rõ họ tên)}\\
\vspace{0.4cm}\\
\textbf{Trịnh Minh Đạt}
\end{tabular}
\end{flushright}
\end{document}